\documentclass{article}

\usepackage{iclr2027_conference,times}
\usepackage[T1]{fontenc}
\usepackage{lmodern}
\usepackage{amsmath,amssymb,amsthm,mathtools}
\usepackage{array,booktabs,microtype,url,xcolor}
\usepackage[colorlinks=true,allcolors=blue]{hyperref}

\newtheorem{theorem}{Theorem}
\newtheorem{lemma}[theorem]{Lemma}
\newtheorem{proposition}[theorem]{Proposition}
\newtheorem{corollary}[theorem]{Corollary}
\theoremstyle{definition}
\newtheorem{definition}[theorem]{Definition}
\theoremstyle{remark}
\newtheorem{remark}[theorem]{Remark}

\newcommand{\Ber}{\operatorname{Ber}}
\newcommand{\E}{\mathbb E}
\newcommand{\1}{\mathbf 1}
\newcommand{\simplex}{\Delta}

\title{No Perfectly Truthful Sequential Calibration Measure}
\author{\textbf{Pahan Dewasurendra} \\ Johns Hopkins University}

\hypersetup{pdftitle={No Perfectly Truthful Sequential Calibration Measure},pdfauthor={Pahan Dewasurendra},pdfkeywords={perfectly, truthful, sequential, calibration, measure, should, assign, sublinear, error},pdfsubject={preprint}}
\begin{document}
\maketitle

\begin{abstract}
A calibration measure should assign sublinear error to correct probabilistic forecasts, linear error to a fixed incorrect forecast, and give a forecaster no incentive to misreport its beliefs. Whether these three requirements can coexist under perfect truthfulness has remained open for sequential prediction. We prove that they cannot. Our main structural result shows that every bounded perfectly truthful measure computed from the outcomes and the forecasts observed along the realized path has a tree-additive Bayes risk. This representation requires no continuity, differentiability, strict truthfulness, or downstream decision guarantee. Although the measure itself may be nonsmooth, one fixed report probability avoids all nondifferentiability points across every horizon. At that report, the measure is forced to be a sum of nonnegative one-step proper losses. Concavity then implies that, under fair independent outcomes, this fixed incorrect report has expected error at most twice the truthful error. Completeness makes the latter sublinear, contradicting soundness. This resolves the perfect-truthfulness question posed by Haghtalab et al. and separates sequential evaluation from recent positive results for batch calibration. The boundary is sharp: if the evaluator can inspect the full forecast tree rather than one realized path, a simple squared-count score satisfies perfect truthfulness, completeness, and product soundness simultaneously.
\end{abstract}

\section{Introduction}

Probabilistic forecasts are useful only when their numerical values can be trusted. Calibration captures a basic form of this requirement: among events assigned probability \(p\), approximately a \(p\) fraction should occur. Calibration has classical roots in forecasting and statistics \citep{brier1950verification,dawid1982well}, supports calibrated learning against arbitrary outcome sequences \citep{foster1998asymptotic}, and is now a standard criterion for uncertainty estimates in machine learning.

A numerical calibration measure affects incentives as well as evaluation. If a forecaster knows the true conditional probability of the next outcome, minimizing the reported calibration error should not reward a lie. \citet{haghtalab2024truthfulness} formalized this requirement in sequential prediction. A measure is perfectly truthful when the strategy that reports every true conditional probability minimizes expected error for every joint distribution of the outcome sequence. They paired truthfulness with two minimal usefulness conditions. Completeness requires sublinear expected error for truthful i.i.d. Bernoulli forecasts. Soundness requires linear error when the same data are evaluated against a fixed incorrect probability.

Known measures expose a sharp tension. Standard calibration errors are complete and sound but can reward strategic forecasts. Subsampled smooth calibration error reduces the incentive gap to a constant multiplicative factor, but it is not perfectly truthful \citep{haghtalab2024truthfulness}. Standard additive proper scoring rules are perfectly truthful, but their one-step noise accumulates to linear error and violates completeness. \citet{qiao2025decision} ruled out perfect truthfulness for measures that also control every downstream decision problem. Their theorem relies on this decision-theoretic condition and does not address the three basic requirements alone. In the batch setting, the answer is positive: \citet{hartline2026perfect} constructed a perfectly truthful, complete, and sound measure using variance additivity. They explicitly retained the sequential case as an open problem.

We close the sequential question with a negative answer. The result holds under exactly the standard information model. At time \(t\), the evaluator eventually sees the outcome \(x_t\) and the forecast \(p_t\) made after observing \(x_{<t}\). It does not see forecasts at histories that did not occur.

\begin{theorem}[Informal]\label{thm:informal}
No bounded sequential calibration measure is simultaneously perfectly truthful, complete, and sound. In fact, deterministic completeness and completeness only for fair i.i.d. outcomes already imply that one fixed incorrect constant forecast has sublinear expected error.
\end{theorem}

The proof is not a lower bound for a particular calibration statistic. It classifies every perfectly truthful realized-path score. A forecast strategy is a probability tree. In the interior, it defines a joint distribution on the \(2^T\) leaves. Perfect truthfulness turns the calibration measure into a proper scoring rule for that leaf distribution. The special information structure is crucial: the score of a realized leaf can inspect only conditional probabilities on its root-to-leaf path.

Our main representation theorem shows that the Bayes risk, or entropy, of every bounded path-local proper score decomposes over the internal nodes of the probability tree. Each node contributes its visitation probability times a concave binary entropy. This is a converse to the familiar construction that sums one-step proper losses. The representation holds even when the original measure is discontinuous and chooses different supporting hyperplanes at nonsmooth points.

The remaining obstacle is precisely that nonsmooth freedom. Each one-dimensional concave node entropy is nondifferentiable at only countably many probabilities. There are countably many nodes over all finite horizons, so one fixed \(\beta\in(1/3,2/3)\), distinct from \(1/2\), avoids every exceptional point. At the constant report \(\beta\), the supporting score is unique and equals the sum of the nodewise tangent losses. Deterministic completeness and boundedness let us normalize every node entropy to vanish at \(0\) and \(1\). The tangent losses are then nonnegative. A short concavity argument gives, at every node, \[ \E_{Y\sim\Ber(1/2)}\ell_{\beta}(Y)\le 2h(1/2). \] Summing with the true history probabilities proves that the wrong constant report costs at most twice the truthful error. Product completeness and product soundness now contradict each other.

The theorem pinpoints the information boundary. Suppose the evaluator receives the entire probability tree. It can compute the tree's reported expectation \(\mu\) of \(\sum_tX_t\) and use \(T^{-1}(\sum_t x_t-\mu)^2\). This score is perfectly truthful. Its truthful risk under i.i.d. Bernoulli data is \(O(1)\), while any fixed incorrect product tree has \(\Omega(T)\) risk. The construction is unavailable from one realized path because \(\mu\) depends on forecasts in unrealized branches.

\paragraph{Contributions.} Our results are:
\begin{itemize}
\item A no-regularity representation theorem for bounded perfectly truthful realized-path scores. Their entropy is a sum of history-dependent binary entropies over the forecast tree.
\item An impossibility theorem resolving the open coexistence of perfect truthfulness, completeness, and soundness in sequential calibration. It removes the decision-theoretic assumption from the previous impossibility result.
\item A quantitative witness. For one fixed \(\beta\ne1/2\), the expected error of reporting \(\beta\) on fair coins is at most twice the truthful expected error at every horizon.
\item A matching information separation. Full-tree access admits a perfectly truthful, deterministically complete, product-complete, and product-sound score.
\end{itemize}

\section{Related work}

\paragraph{Truthful calibration.} \citet{haghtalab2024truthfulness} introduced truthfulness for sequential calibration measures, documented large incentive gaps for standard measures, and constructed subsampled smooth calibration error with constant-factor truthfulness. Their discussion asks whether perfect truthfulness can coexist with completeness and soundness. \citet{qiao2025decision} studied measures that additionally control downstream regret. They constructed subsampled step calibration with strong approximate guarantees and showed that every complete decision-theoretic measure is nontruthful in the unsmoothed setting. Their lower bound uses domination of U-calibration. Our result neither assumes nor proves any downstream decision guarantee. It instead derives a representation from perfect truthfulness itself.

\citet{hartline2026perfect} gave the first perfectly truthful batch calibration measure. Their averaged two-bin error exploits variance additivity for independent batch outcomes. It is complete, sound, continuous, and sample efficient. The same work states the sound, complete, perfectly truthful sequential case as open. A multiclass batch extension was subsequently developed by \citet{lu2025multiclass}. Our theorem explains why batch constructions do not transfer to adaptive sequential reports.

\paragraph{Proper scores and locality.} Proper scoring rules elicit probabilistic beliefs and correspond to supporting hyperplanes of concave entropies \citep{gneiting2007strictly,waghmare2026proper}. \citet{dawid2012local} characterize graph-local proper scores on finite outcome spaces under regularity conditions. \citet{smith2020additive} study strict proper scores that satisfy additivity requirements across marginal and conditional distributions. Prequential forecasting commonly constructs a proper sequence score by summing one-step proper scores \citep{pacchiardi2024probabilistic}. Our representation runs in the converse direction. It assumes only realized-path observability and weak propriety, then derives a history-dependent tree decomposition. The proof uses bounded binary properness rather than differentiability or a graph-clique expansion.

\paragraph{Calibration measures.} Calibration against arbitrary outcome sequences goes back to \citet{foster1998asymptotic}, with later work studying smooth calibration and strategic hedging \citep{kakade2008deterministic,foster2021hedging}. The systematic comparison of calibration measures includes distance, continuity, statistical estimation, and decision consequences \citep{blasiok2023unifying}. Our result concerns a different axis: exact incentive compatibility when forecasts can adapt to observed outcomes.

\begin{table}[t]
\centering
\footnotesize \setlength{\tabcolsep}{3pt}
\begin{tabular}{@{}lcccc@{}}
\toprule
Result & Setting & Complete & Sound & Perfect truth \\
\midrule
\citet{haghtalab2024truthfulness}, SSCE & sequential & yes & yes & no, constant factor \\
\citet{qiao2025decision} & seq., decision-theoretic & yes & yes & impossible \\
\citet{hartline2026perfect}, ATB & independent batch & yes & yes & yes \\
This work & sequential path & yes & yes & impossible \\
This work, full-tree score & full forecast tree & yes & product & yes \\
\bottomrule
\end{tabular}
\caption{The information model separates the positive batch and full-tree results from the sequential impossibility.}
\label{tab:comparison}
\end{table}

\section{Setup}

Write \(\{0,1\}^{<T}=\bigcup_{t=0}^{T-1}\{0,1\}^t\), and let \(x_{<t}=(x_1,\ldots,x_{t-1})\). A deterministic forecaster is a binary probability tree \(A=(a_u)_{u\in\{0,1\}^{<T}}\), where \(a_u\in[0,1]\) is the prediction after history \(u\). On outcome sequence \(x\in\{0,1\}^T\), its realized prediction vector is \[ A(x)=(a_{x_{<1}},\ldots,a_{x_{<T}}). \] A calibration measure is a family \(M=(M_T)_{T\ge1}\) with \(M_T:\{0,1\}^T\times[0,1]^T\to[0,T]\). Its expected error under outcome distribution \(P\) and strategy \(A\) is \[ \operatorname{err}_{M_T}(P,A)=\E_{X\sim P}M_T(X,A(X)). \]

For a full-support \(P\in\simplex(\{0,1\}^T)\), its truthful tree has node reports \[ q_u=P(X_{|u|+1}=1\mid X_{1:|u|}=u). \] At zero-probability histories the truthful report can be chosen arbitrarily.

\begin{definition}[Truthfulness, completeness, and soundness]\label{def:properties}
The measure \(M\) is perfectly truthful if the truthful tree minimizes expected error for every \(T\) and every \(P\). It is deterministically complete if \(M_T(x,x)=0\) for all \(T,x\). It is product-complete if, for every fixed \(\alpha\in[0,1]\), \[ \E_{X\sim\Ber(\alpha)^{\otimes T}}M_T(X,\alpha\1)=o_\alpha(T). \] It is product-sound if, for every fixed \(\alpha\ne\beta\), \[ \E_{X\sim\Ber(\alpha)^{\otimes T}}M_T(X,\beta\1)=\Omega_{\alpha,\beta}(T). \]
\end{definition}

Deterministic and product completeness are exactly the two completeness clauses of \citet{haghtalab2024truthfulness}. Product soundness is one of their two soundness clauses. Their other soundness clause requires linear error on the pointwise opposite prediction. The later definition of \citet{qiao2025decision} adds sublinear error on every empirically calibrated pair. Our impossibility does not need either addition.

Every interior report tree \(A\), with \(a_u\in(0,1)\), induces a unique full-support distribution \(Q_A\) on the leaves: \[ Q_A(x)=\prod_{t=1}^T a_{x_{<t}}^{x_t}(1-a_{x_{<t}})^{1-x_t}. \] Conversely, the truthful tree of \(Q_A\) is \(A\). Define the induced leaf score and its Bayes risk by \[ S_T(Q,x)=M_T(x,A_Q(x)),\qquad H_T(P)=\E_{X\sim P}S_T(P,X). \tag{1} \] Perfect truthfulness says that \(S_T\) is a proper scoring rule: \[ H_T(P)\le \E_{X\sim P}S_T(Q,X) \tag{2} \] for all full-support \(P,Q\). It follows that \(H_T\) is concave. Crucially, \(S_T(Q,x)\) depends only on the conditional probabilities of \(Q\) along the realized path to \(x\).

\section{The tree-additive representation}

The representation theorem is stated for binary trees because that is the calibration framework of Definition~\ref{def:properties}. For a history \(u\), write \(P(u)=P(X_{1:|u|}=u)\).

\begin{theorem}[Tree-additive entropy]\label{thm:factorization}
Fix \(T\) and a bounded perfectly truthful realized-path score. There are bounded concave functions \(h_{T,u}:[0,1]\to\mathbb R\), one for every \(u\in\{0,1\}^{<T}\), such that every full-support distribution \(P\) satisfies \[ H_T(P)=\sum_{u\in\{0,1\}^{<T}}P(u)h_{T,u}(q_u). \tag{3} \] If the score is nonnegative and deterministically complete, the functions can be chosen so that \[ h_{T,u}(0)=h_{T,u}(1)=0,\qquad h_{T,u}(q)\ge0. \tag{4} \] At any report distribution \(Q\) where every \(h_{T,u}\) is differentiable at the corresponding conditional \(q_u\), the actual realized score is \[ S_T(Q,x)=\sum_{u\prec x}\left[h_{T,u}(q_u)+(x_{|u|+1}-q_u)h_{T,u}'(q_u)\right]. \tag{5} \]
\end{theorem}

Equation~(3) says that the expected truthful score is assembled from binary uncertainty at each visited history. Equation~(5) is stronger at smooth reports. It identifies the actual measure, not only its expectation.

The proof starts with a binary rigidity fact. It is useful because the conditional score below a branch can depend on the reported probability above that branch.

\begin{lemma}[Shared-side uniqueness]\label{lem:shared}
Let \((\ell_0,\ell_1)\) and \((\widetilde\ell_0,\ell_1)\) be bounded proper binary losses on \((0,1)\). Then \(\ell_0-\widetilde\ell_0\) is constant. The symmetric statement holds when the zero-outcome component is shared.
\end{lemma}

For \(0<p<q<1\), properness at true parameters \(p\) and \(q\) traps the increment of \(\ell_0\): \[ \frac{p}{1-p}\bigl(\ell_1(p)-\ell_1(q)\bigr) \le \ell_0(q)-\ell_0(p) \le \frac{q}{1-q}\bigl(\ell_1(p)-\ell_1(q)\bigr). \tag{6} \] The same interval traps the increment of \(\widetilde\ell_0\). On a refining partition, the sum of the interval widths tends to zero because \(\ell_1\) is bounded and monotone. Thus the two increments agree. Appendix~\ref{app:shared} gives the details.

To see how Lemma~\ref{lem:shared} yields Theorem~\ref{thm:factorization}, split a leaf distribution at the root as \[ P=((1-a)P_0,aP_1). \] Path locality writes the score on root branch \(b\) as \(S_b(p,Q_b,x)\), with no dependence on \(Q_{1-b}\). For fixed root report \(p\), this is a proper score for the conditional leaf distribution in branch \(b\). Let its entropy be \(K_b^p(R)\). Varying only the root report shows that \[ p\mapsto\bigl(K_0^p(R_0),K_1^p(R_1)\bigr) \] is a proper binary loss for every \(R_0,R_1\). Holding one conditional fixed and applying Lemma~\ref{lem:shared} gives \[ K_b^p(R)=a_b(p)+K_b(R). \tag{7} \] The pair \((a_0,a_1)\) is proper at the root, and each \(K_b\) is the entropy of a bounded path-local proper score on a depth-\((T-1)\) tree. Recursion proves Equation~(3).

The endpoint normalization needs no continuity. If \(P\) approaches a deterministic leaf \(x\), perfect truthfulness allows comparison with the strategy hard-coded to predict \(x\). It loses zero on \(x\) and at most \(T\) elsewhere, so \[ 0\le H_T(P)\le T(1-P(x))\longrightarrow0. \tag{8} \] The endpoint values of the node entropies therefore sum to zero on every root-to-leaf path. Affine node potentials set both endpoints of every \(h_{T,u}\) to zero while telescoping out of Equation~(3). Concavity then gives nonnegativity. Finally, the score vector in Equation~(2) is a supporting hyperplane of \(H_T\). When all node entropies are differentiable, this hyperplane is unique. Differentiating the perspective term \(P(u)h_{T,u}(q_u)\) gives Equation~(5).

\section{Perfect truthfulness contradicts soundness}

The local comparison behind the impossibility is elementary.

\begin{lemma}[Tangent comparison]\label{lem:tangent}
Let \(h:[0,1]\to\mathbb R_+\) be concave with \(h(0)=h(1)=0\), and suppose \(h\) is differentiable at \(p\in(0,1)\). Define \[ \ell_p(b)=h(p)+(b-p)h'(p),\qquad b\in\{0,1\}. \] For every \(q\in(0,1)\), \[ (1-q)\ell_p(0)+q\ell_p(1)\le C_{p,q}h(q), \tag{9} \] where \[ C_{p,q}=\max\left\{\frac q p,\frac{1-q}{1-p}\right\} \max\left\{\frac p q,\frac{1-p}{1-q}\right\}. \]
\end{lemma}

The tangent line lies above the zero endpoint values, hence \(\ell_p(0),\ell_p(1)\ge0\). Reweighting these two nonnegative numbers gives \[ (1-q)\ell_p(0)+q\ell_p(1) \le \max\left\{\frac q p,\frac{1-q}{1-p}\right\}h(p). \] Concavity between \(p,q\), and the appropriate endpoint gives the second factor in Equation~(9).

\begin{theorem}[Sequential impossibility]\label{thm:impossibility}
No calibration measure \(M_T:\{0,1\}^T\times[0,1]^T\to[0,T]\) is simultaneously perfectly truthful, deterministically complete, product-complete, and product-sound.
\end{theorem}

\begin{proof}
Apply Theorem~\ref{thm:factorization} at every horizon. Each concave node entropy is nondifferentiable at at most countably many points. The union of these exceptional sets over all nodes and horizons is countable. Choose one fixed \[ \beta\in(1/3,2/3)\setminus\{1/2\} \] outside this union, and set \(\alpha=1/2\).

At the product report \(Q_\beta=\Ber(\beta)^{\otimes T}\), every node conditional is \(\beta\). Equation~(5) and Lemma~\ref{lem:tangent} give
\begin{align*}
\E_{X\sim\Ber(1/2)^{\otimes T}}M_T(X,\beta\1)
&\le C_{\beta,1/2}\sum_u P_{1/2}(u)h_{T,u}(1/2)\\
&=C_{\beta,1/2}H_T(P_{1/2})\\
&=C_{\beta,1/2}\E_{X\sim\Ber(1/2)^{\otimes T}}M_T(X,\tfrac12\1).
\end{align*}
For \(\beta\in(1/3,2/3)\), the constant is \[ C_{\beta,1/2}=\max\left\{\frac{\beta}{1-\beta},\frac{1-\beta}{\beta}\right\}<2. \] Product completeness makes the last expectation \(o(T)\). Product soundness requires the first expectation to be \(\Omega(T)\), a contradiction.
\end{proof}

The same history law \(P_{1/2}(u)\) weights both sides of the nodewise comparison. No likelihood-ratio comparison between the truthful and incorrect product distributions is used. This point rules out an apparent escape in which a measure concentrates penalties on histories that are rare under the reported model.

\begin{corollary}[Quantitative failure]\label{cor:quantitative}
For every bounded, perfectly truthful, deterministically complete family \(M\), there is a fixed \(\beta\in(1/3,2/3)\setminus\{1/2\}\) such that, for every \(T\), \[ \E_{P_{1/2}}M_T(X,\beta\1)\le2\E_{P_{1/2}}M_T(X,\tfrac12\1). \]
\end{corollary}

Thus the obstruction is stronger than the failure of a particular asymptotic rate. A fixed wrong report cannot even be separated from truth by more than a universal factor.

\section{Full-tree access restores possibility}

The standard evaluator observes \(A(X)\), not the reports at histories outside the realized path. We now show that this restriction is exactly where the impossibility enters.

Suppose an evaluator receives the full probability tree \(A\). Let \(Q_A\) be its induced leaf distribution and define \[ \mu(A)=\E_{Y\sim Q_A}\sum_{t=1}^T Y_t, \qquad M_T^{\mathrm{tree}}(x,A)=\frac1T\left(\sum_{t=1}^T x_t-\mu(A)\right)^2. \tag{10} \]

\begin{proposition}[Full-tree possibility]\label{prop:fulltree}
The score in Equation~(10) takes values in \([0,T]\), is perfectly truthful for every joint outcome distribution, is zero on deterministic truth, and is product-complete and product-sound.
\end{proposition}

\begin{proof}
For true distribution \(P\), write \(Z=\sum_tX_t\). Then \[ \E_P M_T^{\mathrm{tree}}(X,A) =\frac{\operatorname{Var}_P(Z)}{T} +\frac{(\E_PZ-\mu(A))^2}{T}. \] The truthful tree induces \(P\) and reports \(\mu(A_P)=\E_PZ\), so it is optimal. For deterministic \(P=\delta_x\), its score is zero. Under \(P_\alpha=\Ber(\alpha)^{\otimes T}\), the truthful expected score is \(\alpha(1-\alpha)\). The constant tree \(A\equiv\beta\) has expected score \[ \alpha(1-\alpha)+T(\alpha-\beta)^2. \] This is linear for every fixed \(\alpha\ne\beta\).
\end{proof}

The score only elicits the expected total count, so truth need not be its unique minimizer. Perfect truthfulness in the calibration definition also does not require uniqueness. Computing \(\mu(A)\) requires integrating over unrealized branches. One path of scalar conditional forecasts contains insufficient information.

\section{Discussion}

Perfect truthfulness in sequential calibration is impossible for a structural reason. A measure that sees only realized forecasts becomes a path-local proper score. Its uncertainty must accumulate through conditional binary entropies. Completeness asks this accumulated truthful uncertainty to be sublinear. Concavity then prevents the same node losses from assigning linear error to a fixed incorrect probability.

The proof leaves several directions open. Approximate truthfulness remains possible and quantitatively important. Existing constant-factor constructions show that an exact-to-approximate transition changes the problem sharply. It would be useful to determine the best universal factor under the original completeness and soundness axioms. The full-tree result also suggests studying limited counterfactual access, such as a bounded number of queried unrealized histories or auxiliary forecasts of aggregate outcomes. Finally, a multiclass extension should replace binary node entropies by concave functions on a simplex. The tangent comparison extends directly, while a no-regularity shared-component factorization for categorical nodes requires additional care.

\bibliography{references}
\bibliographystyle{iclr2027_conference}

\appendix

\section{Proper-score preliminaries}\label{app:proper}

We record the finite-dimensional facts used in the proof. Let \(\mathcal Y\) be finite. A loss-oriented score \(S(Q,y)\) is proper if \[ \E_{Y\sim P}S(P,Y)\le \E_{Y\sim P}S(Q,Y) \] for all \(P,Q\in\simplex(\mathcal Y)\). Its entropy is \(H(P)=\E_PS(P,Y)\). Since \[ H(P)=\inf_Q\E_PS(Q,Y), \] it is an infimum of affine functions of \(P\) and is therefore concave. For each \(Q\), the vector \(S(Q,\cdot)\) defines an affine majorant of \(H\): \[ H(P)\le \sum_yP(y)S(Q,y), \qquad H(Q)=\sum_yQ(y)S(Q,y). \tag{11} \] Conversely, a choice of supporting affine majorant at every \(Q\) defines a proper score. This is the finite-outcome Savage representation \citep{gneiting2007strictly,dawid2012local}.

If \(H\) is differentiable at an interior \(Q\), the supporting affine function in Equation~(11) is unique. To see the normalization explicitly, extend \(H\) homogeneously to the positive cone by \[ \overline H(z)=\left(\sum_yz_y\right)H\left(\frac{z}{\sum_yz_y}\right). \] The score vector is a supergradient of the concave one-homogeneous function \(\overline H\). At differentiability points it equals \(\nabla\overline H(Q)\).

\section{Proof of shared-side uniqueness}\label{app:shared}

We prove Lemma~\ref{lem:shared}. A bounded pair \((\ell_0,\ell_1)\) is proper when, for all \(p,q\in(0,1)\), \[ (1-q)\ell_0(q)+q\ell_1(q) \le(1-q)\ell_0(p)+q\ell_1(p). \tag{12} \] Fix \(0<p<q<1\). Apply Equation~(12) first with true parameter \(p\) and report \(q\), then with true parameter \(q\) and report \(p\). With \[ \Delta_0=\ell_0(q)-\ell_0(p), \qquad \Delta_1=\ell_1(p)-\ell_1(q), \] the two inequalities become \[ \frac{p}{1-p}\Delta_1\le\Delta_0\le\frac{q}{1-q}\Delta_1. \tag{13} \] The interval is nonempty only if \(\Delta_1\ge0\), and then \(\Delta_0\ge0\). Thus \(\ell_0\) is nondecreasing and \(\ell_1\) is nonincreasing.

Let \((\widetilde\ell_0,\ell_1)\) also be proper. Its increment \(\widetilde\Delta_0\) lies in the same interval in Equation~(13). Set \(g(r)=r/(1-r)\). On any compact interval \([a,b]\subset(0,1)\), take a partition \(a=x_0<\cdots<x_n=b\). Then
\begin{align*}
&\left|[\ell_0(b)-\ell_0(a)]-[\widetilde\ell_0(b)-\widetilde\ell_0(a)]\right|\\
&\quad\le\sum_{i=1}^n[g(x_i)-g(x_{i-1})][\ell_1(x_{i-1})-\ell_1(x_i)]\\
&\quad\le\max_i[g(x_i)-g(x_{i-1})][\ell_1(a)-\ell_1(b)].
\end{align*}
The last expression tends to zero along partitions whose mesh tends to zero. Therefore the two increments agree on every compact interval, and \(\ell_0-\widetilde\ell_0\) is constant. Swapping outcomes proves the symmetric statement.

\section{Proof of the entropy factorization}\label{app:factorization}

We prove Equation~(3) of Theorem~\ref{thm:factorization} by induction on \(T\). The case \(T=1\) is the ordinary binary entropy of a bounded proper score.

For \(T\ge2\), split the leaf space into \(\mathcal Y_0=\{x:x_1=0\}\) and \(\mathcal Y_1=\{x:x_1=1\}\). Parameterize a full-support report distribution as \[ Q=((1-p)Q_0,pQ_1). \] For \(x\in\mathcal Y_b\), realized-path locality implies \[ S_T(Q,x)=S_b(p,Q_b,x), \tag{14} \] with no dependence on \(Q_{1-b}\).

Fix \(p\in(0,1)\). For every true conditional \(R_b\), global propriety with root probability \(p\) and all other reports fixed at truth implies \[ \E_{R_b}S_b(p,R_b,Y)\le\E_{R_b}S_b(p,Q_b,Y). \] Thus \(S_b(p,\cdot,\cdot)\) is a bounded proper score on \(\mathcal Y_b\). Denote its entropy by \[ K_b^p(R_b)=\E_{Y\sim R_b}S_b(p,R_b,Y). \]

Now fix true conditionals \(R_0,R_1\) and vary only the root report. Global propriety gives, for every true root probability \(a\in(0,1)\), \[ (1-a)K_0^a(R_0)+aK_1^a(R_1) \le(1-a)K_0^p(R_0)+aK_1^p(R_1). \tag{15} \] Hence \(p\mapsto(K_0^p(R_0),K_1^p(R_1))\) is a bounded proper binary loss for every pair \(R_0,R_1\).

Fix a reference \(R_0^\circ\). Holding \(R_1\) fixed in Equation~(15), Lemma~\ref{lem:shared} shows that \[ K_0^p(R_0)-K_0^p(R_0^\circ) \] is independent of \(p\). Thus \(K_0^p(R_0)=a_0(p)+K_0(R_0)\). The symmetric argument gives \(K_1^p(R_1)=a_1(p)+K_1(R_1)\). Substituting these expressions into Equation~(15) shows that \((a_0,a_1)\) is a proper binary loss because the \(K_b(R_b)\) terms are independent of the root report. Its entropy \[ h_{T,\varnothing}(a)=(1-a)a_0(a)+aa_1(a) \] is concave.

Choose any \(p_0\in(0,1)\). The branch score \[ \widetilde S_b(Q_b,y)=S_b(p_0,Q_b,y)-a_b(p_0) \] is bounded, proper, and path-local, with entropy \(K_b\). The induction hypothesis decomposes \(K_b\) over the internal nodes below branch \(b\). Finally,
\begin{align*}
H_T(P)
&=(1-a)K_0^a(P_0)+aK_1^a(P_1)\\
&=h_{T,\varnothing}(a)+(1-a)K_0(P_0)+aK_1(P_1),
\end{align*}
which yields Equation~(3).

\section{Boundary normalization and score identification}\label{app:normalization}

We first prove Equation~(8) carefully. For a leaf \(x\in\{0,1\}^T\), define a valid hard-coded strategy \(A^x\) by \(A^x(u)=x_{|u|+1}\) at every history of the appropriate depth. On outcome \(x\), its realized prediction vector equals \(x\), so deterministic completeness gives \(M_T(x,A^x(x))=0\). Since every loss is at most \(T\), perfect truthfulness gives \[ 0\le H_T(P)\le\E_PM_T(X,A^x(X))\le T(1-P(x)). \tag{16} \]

Each bounded concave node entropy in Equation~(3) has finite one-sided limits at \(0\) and \(1\). Write these limits as \(e_{u,0}\) and \(e_{u,1}\). Approach a deterministic leaf \(x\) through interior trees whose conditional probability at every node on the path tends to the corresponding bit of \(x\). Terms at off-path nodes vanish because their visitation probabilities tend to zero. Equations~(3) and (16) imply \[ \sum_{u\prec x}e_{u,x_{|u|+1}}=0 \tag{17} \] for every leaf \(x\).

Set \(c_\varnothing=0\) and recursively define \(c_{ub}=c_u-e_{u,b}\). Equation~(17) gives \(c_x=0\) at every leaf. Replace \[ h_{T,u}(q) \quad\text{by}\quad \widetilde h_{T,u}(q)=h_{T,u}(q)+(1-q)c_{u0}+qc_{u1}-c_u. \tag{18} \] The change in the entropy sum is
\begin{align*}
\sum_uP(u)[(1-q_u)c_{u0}+q_uc_{u1}-c_u]
&=\sum_u[P(u0)c_{u0}+P(u1)c_{u1}-P(u)c_u]\\
&=\sum_xP(x)c_x-c_\varnothing=0.
\end{align*}
Both endpoint values in Equation~(18) are zero. The transformed functions remain concave, so they are nonnegative. We drop the tilde below.

Suppose every \(h_{T,u}\) is differentiable at the conditional \(q_u\) of an interior report \(Q\). In leaf-probability coordinates, the node contribution is the perspective \[ [Q(u0)+Q(u1)]h_{T,u}\left(\frac{Q(u1)}{Q(u0)+Q(u1)}\right). \] Its derivative with respect to the probability of a descendant leaf in child \(b\) is \[ h_{T,u}(q_u)+(b-q_u)h_{T,u}'(q_u). \] It has derivative zero for leaves outside the subtree. Summing over nodes gives the right side of Equation~(5). By Appendix~\ref{app:proper}, this derivative is the unique supporting score at \(Q\), so it equals the actual \(S_T(Q,x)\).

\section{Proof of the tangent comparison}\label{app:tangent}

Let \[ \ell_p(0)=h(p)-ph'(p), \qquad \ell_p(1)=h(p)+(1-p)h'(p). \] The tangent inequality for a concave function gives \[ 0=h(0)\le h(p)-ph'(p), \qquad 0=h(1)\le h(p)+(1-p)h'(p). \] Therefore both losses are nonnegative, and
\begin{align*}
(1-q)\ell_p(0)+q\ell_p(1)
&\le\max\left\{\frac{1-q}{1-p},\frac q p\right\}
[(1-p)\ell_p(0)+p\ell_p(1)]\\
&=\max\left\{\frac{1-q}{1-p},\frac q p\right\}h(p).
\tag{19}
\end{align*}
If \(p\le q\), write \(q=\lambda p+(1-\lambda)1\), where \(\lambda=(1-q)/(1-p)\). Concavity gives \[ h(q)\ge\frac{1-q}{1-p}h(p). \] If \(p\ge q\), write \(q=\lambda p+(1-\lambda)0\), where \(\lambda=q/p\), to obtain \(h(q)\ge(q/p)h(p)\). In both cases, \[ h(p)\le\max\left\{\frac{1-p}{1-q},\frac p q\right\}h(q). \tag{20} \] Equations~(19) and (20) prove Lemma~\ref{lem:tangent}.

\section{A two-step illustration}\label{app:two-step}

For \(T=2\), write a report tree as \((a,b_0,b_1)\), where \(a\) is the first forecast and \(b_i\) is the second forecast after observing \(X_1=i\). A realized-path score has the form \[ S_{00}(a,b_0),\quad S_{01}(a,b_0),\quad S_{10}(a,b_1),\quad S_{11}(a,b_1). \] Theorem~\ref{thm:factorization} says that its entropy must be \[ H(a,c_0,c_1)=h_\varnothing(a)+(1-a)h_0(c_0)+ah_1(c_1). \] At a differentiability point, the four losses are
\begin{align*}
S_{0j}&=h_\varnothing(a)-ah_\varnothing'(a)+h_0(b_0)+(j-b_0)h_0'(b_0),\\
S_{1j}&=h_\varnothing(a)+(1-a)h_\varnothing'(a)+h_1(b_1)+(j-b_1)h_1'(b_1).
\end{align*}
The score may use a different binary entropy after each first-round history. It cannot include an entropy interaction between \(b_0\) and \(b_1\), because the realized score sees only one of them. This small case displays the general causal factorization.

\section{Why a product-family score is not enough}\label{app:product}

If forecasts were restricted to one constant parameter \(p\), the score \[ \frac1T\left(\sum_{t=1}^T(x_t-p)\right)^2 =T(\overline x-p)^2 \] would be proper for the Bernoulli product family. Its truthful expected value is \(\alpha(1-\alpha)\), while a fixed wrong report has an additional \(T(\alpha-p)^2\) cost. This does not contradict Theorem~\ref{thm:impossibility}. For a general sequential distribution, the forecasts \(p_t\) can adapt to earlier residuals. Squaring the terminal sum creates cross terms through which a strategic forecaster can hedge past noise. The full-tree score in Equation~(10) avoids this problem by using the ex ante expectation of the total count under the entire reported tree.

\end{document}
